\documentclass[12pt]{article}
\usepackage[utf8]{inputenc}
\usepackage[T1]{fontenc}
\usepackage{hyperref}
\title{Kiyombe Language Resources for K-12 Curriculum Development}
\author{Compiled by Claude Sonnet 5.5 (claude-sonnet-5-5) for Andr\'as Kornai}
\date{}
\begin{document}
\maketitle

\section{Introduction}
Kiyombe (Yombe; ISO 639-3 \texttt{yom}, Glottocode yomb1244) is a Kikongo variety listed separately under its own ISO code. It is spoken in the Democratic Republic of the Congo, Angola and Congo. Web searches found very little Kiyombe-specific educational material. Most of what exists concerns Kikongo broadly. This report lists only what appeared in search results and says where a resource is not specific to Kiyombe.

\section{Primary Educational Resources}
No Kiyombe primers, textbooks or school programmes were found in the searches. Nothing indicates that the language is taught in schools.

\section{Linguistic and Reference Materials}
\begin{itemize}
\item A French--Kiyombe lexicon by Biyoko Mabua, described as containing many example sentences: \url{https://e.lavoisier.fr/extract/523164}
\item W.\ Holman Bentley, \emph{Dictionary and grammar of the Kongo language} (Baptist Missionary Society, 1887--1895). It describes the San Salvador variety, not Kiyombe specifically, and is useful only as a related reference: \url{https://catalog.hathitrust.org/Record/008584146}
\item General background on the Kongo language cluster: \url{https://en.wikipedia.org/wiki/Kongo_language}
\end{itemize}

\section{Religious and Community Resources}
\begin{itemize}
\item Open Kiyombe Contemporary New Testament and Psalms (Biblica, 2002; released under CC BY-SA 4.0), as PDF: \url{https://eBible.org/pdf/yom/yom_all.pdf}
\item The same text as EPUB: \url{https://www.ebible.org/epub/yom.epub}
\end{itemize}

\section{Digital and Audio Resources}
Scripture Earth lists audio and video for Kikongo (code kng), including The JESUS Film. The search did not confirm that these are in the Kiyombe variety: \url{https://www.scriptureearth.org/00eng.php?iso=kng}

\section{Curriculum Development Context}
The only Kiyombe-specific items found are one lexicon and one Bible translation. A native-language K-12 curriculum would therefore need new materials. Kikongo materials may be adaptable, but they differ by variety, and Kikongo has no single standard orthography.
\end{document}